\section*{Chapter \ref{ch:b2:muscle-movement} --- Muscle Contraction and Motor Function}

\begin{solution}{exo:b2:muscle-movement:1}
Thin filaments slide along thick ones, pulled by myosin heads;
neither filament shortens. The I band and the H zone narrow and the
Z discs approach; the A band (the thick filaments) keeps its length.
\end{solution}

\begin{solution}{exo:b2:muscle-movement:2}
(1) ATP bound, head detached; ATP hydrolysed, head cocked. (2) Head
binds actin. (3) Phosphate released, power stroke, ADP released.
(4) New ATP binds and the head detaches. ATP is needed to
\emph{release} the head: without it every head stays bound and the
muscle locks --- rigor.
\end{solution}

\begin{solution}{exo:b2:muscle-movement:3}
Nerve impulse; release of acetylcholine and the end-plate potential
(\qty{0.6}{ms}); muscle \omterm{prop:b2:neurons-synapses:ap}{action potential} over the fibre and down the
\omterm{def:b2:cell-differentiation:fibre}{T tubules} (\qty{1}{ms}); calcium released from the \omterm{def:b2:cell-differentiation:fibre}{sarcoplasmic
reticulum} (\qty{1}{ms}); calcium binds troponin, tropomyosin moves,
heads attach (a few milliseconds); force rises over tens of
milliseconds.
\end{solution}

\begin{solution}{exo:b2:muscle-movement:4}
\omterm{def:b2:muscle-movement:motorunit}{Motor unit}: one motor \omterm{def:b2:neurons-synapses:neuron}{neuron} and all the fibres it innervates.
\omterm{def:b2:muscle-movement:motorunit}{Twitch}: the brief contraction from a single impulse. \omterm{def:b2:muscle-movement:motorunit}{Tetanus}: the
fused, sustained contraction from impulses arriving faster than the
\omterm{def:b2:muscle-movement:motorunit}{twitch} decays. Force is graded by the firing rate of each unit and
by the number of units recruited.
\end{solution}

\begin{solution}{exo:b2:muscle-movement:5}
Each \omterm{def:b2:cell-differentiation:fibre}{half-sarcomere} slides \qty{0.2}{\micro m} in \qty{50}{ms}:
\qty{4}{\micro m/s}; that is 20 strokes of \qty{10}{nm} in
\qty{50}{ms}, 400 a second if one head were attached throughout.
\end{solution}

\begin{solution}{exo:b2:muscle-movement:6}
\qty{2}{cm} in \qty{0.1}{s}: \qty{0.2}{m/s}, 2 lengths a second; each
\omterm{def:b2:cell-differentiation:fibre}{sarcomere} shortens $2/\num{40000}$ cm $= \qty{0.5}{\micro m}$ in
\qty{0.1}{s}, \qty{5}{\micro m/s}.
\end{solution}

\begin{solution}{exo:b2:muscle-movement:7}
$F/F_0 = 0.3125/(v/v_{\max} + 0.25) - 0.25$: at 1, 3 and 6 lengths a
second ($0.1$, $0.3$, $0.6$ of $v_{\max}$): $0.64$, $0.32$, $0.12$.
Power $Fv$ (in $F_0\times$ lengths per second): $0.64$, $0.95$, $0.71$
--- greatest near 3 lengths a second, about a third of $v_{\max}$.
\end{solution}

\begin{solution}{exo:b2:muscle-movement:8}
$80\times 30 = \qty{2400}{N}$; at the foot $2400\times 5/45 =
\qty{270}{N}$.
\end{solution}

\begin{solution}{exo:b2:muscle-movement:9}
At \qty{10}{Hz} the impulses come every \qty{100}{ms}, after the
calcium of the previous one is gone: separate \omterm{def:b2:muscle-movement:motorunit}{twitches}, each
starting from a partly relaxed state, giving a ripple. At
\qty{50}{Hz} they come every \qty{20}{ms}, before the calcium is
pumped away: the calcium stays high and the force fuses. The
\omterm{def:b2:muscle-movement:motorunit}{tetanus} exceeds the \omterm{def:b2:muscle-movement:motorunit}{twitch} because a single calcium pulse is too
brief for the cross-bridges to stretch the series elastic elements
and develop full force; sustained calcium lets them.
\end{solution}

\begin{solution}{exo:b2:muscle-movement:10}
Work \qty{20}{kJ}, ATP energy \qty{80}{kJ}, $1.6$ mol of ATP. First
\qty{5}{s}: $0.8$ mol of \omterm{prop:b2:muscle-movement:energy}{creatine phosphate}. Remaining $0.8$ mol of ATP
by glycolysis: $0.4$ mol of glucose, $0.8$ mol of lactate --- in
\qty{40}{L}, \qty{20}{mmol/L}.
\end{solution}

\begin{solution}{exo:b2:muscle-movement:11}
Blocked release channel: no calcium, no contraction --- paralysis.
Blocked pump: calcium stays up, the muscle cannot relax --- a
contracture. Insensitive troponin: tropomyosin never moves, no
contraction. No calcium in the \omterm{def:b2:blood-circulation:blood}{blood}: skeletal muscle contracts
normally (its calcium is internal); the \omterm{def:b2:heart:anatomy}{heart}, which needs
external calcium to trigger its release, stops.
\end{solution}

\begin{solution}{exo:b2:muscle-movement:12}
The head converts up to half of an ATP's free energy into work;
the rest is heat, and the losses of the pumps, of the cross-bridges
that attach without pulling at speed, and of the elastic elements
bring the whole to a quarter --- the efficiency of a good petrol
engine. Two ways of grading force allow both fine control at low
effort (small units added one at a time, each firing faster or
slower) and a large range (a thousandfold, from one small unit at
\omterm{def:b2:muscle-movement:motorunit}{twitch} rate to every unit in \omterm{def:b2:muscle-movement:motorunit}{tetanus}); a single switch would give
one force or none.
\end{solution}

\begin{solution}{pb:b2:muscle-movement:1}
\textbf{1.} $v = \sqrt{2\times 9.81\times 0.4} = \qty{2.8}{m/s}$;
$\tfrac12\times 70\times 2.8^{2} = \qty{275}{J}$.
\textbf{2.} Work $= 275 + 70\times 9.81\times 0.4 = \qty{550}{J}$ over
\qty{0.4}{m}: \qty{1370}{N}, twice the weight.
\textbf{3.} $F_0 = 500\times 30 = \qty{15000}{N}$; at the foot,
\qty{1500}{N}.
\textbf{4.} \qty{4}{cm} in \qty{0.3}{s}: \qty{0.13}{m/s}, $1.1$ lengths a
second, $0.14\,v_{\max}$.
\textbf{5.} $F/F_0 = 0.3125/(0.14 + 0.25) - 0.25 = 0.55$: \qty{8300}{N}
in the muscle, \qty{830}{N} at the foot --- well short of the
\qty{1370}{N} needed. Muscle shortening alone cannot make this jump.
The counter-movement supplies the rest: as the jumper dips, the
stretched tendons store elastic energy, and the muscle, contracting
nearly isometrically at high force, loads them; the tendons then
recoil faster than any fibre can shorten.
\textbf{6.} $550/0.3 = \qty{1.8}{kW}$; \qty{92}{W/kg} of extensor.
\textbf{7.} $4/\num{48000}$ cm $= \qty{0.83}{\micro m}$ per \omterm{def:b2:cell-differentiation:fibre}{sarcomere},
\qty{0.42}{\micro m} per half: 42 strokes of \qty{10}{nm}.
\textbf{8.} \qty{0.42}{\micro m} in \qty{0.3}{s}: \qty{1.4}{\micro m/s};
140 strokes a second if continuously attached.
\textbf{9.} $500\times 6\times 10^{10}\times \num{48000}\times 300 =
4.3\times 10^{20}$ heads.
\textbf{10.} $0.55\times \num{15000}\times 0.04 = \qty{330}{J}$; per
stroke $0.4\times 8\times 10^{-20} = \qty{3.2e-20}{J}$; $1.0\times
10^{22}$ strokes.
\textbf{11.} $1.0\times 10^{22}/4.3\times 10^{20} = 24$ strokes per
head, of the 42 available: about \qty{60}{\%}. The reserve is
small --- the muscle is working near its limit, as question 5
found.
\textbf{12.} $1.0\times 10^{22}/6.0\times 10^{23} = 0.017$ mol,
\qty{8.7}{g}; its energy $0.017\times 50 = \qty{860}{J}$, of which
\qty{330}{J} became work: \qty{39}{\%}.
\textbf{13.} Faster than one impulse per \qty{30}{ms}: above
\qty{33}{Hz}; about 10 impulses in \qty{0.3}{s}.
\textbf{14.} $9.9\times 10^{-6}\times 5\times 10^{-10}\times 6\times
10^{23} = 3\times 10^{9}$ free ions; $1.5\times 10^{9}$ ATP.
\textbf{15.} Fibres: $500/(\pi\times(3\times 10^{-3})^{2}) = 1.8\times
10^{7}$; cross-bridge ATP per fibre per impulse $1.0\times
10^{22}/(1.8\times 10^{7}\times 10) = 6\times 10^{13}$, forty thousand
times the pump estimate. The free calcium is a small remainder: the
reticulum releases of the order of \qty{0.1}{mmol/L}, a hundred
times the free rise, almost all of it bound at once by troponin and
the buffers --- and all of it must be pumped back.
\textbf{16.} For a maximal jump, nearly every unit within a few tens
of milliseconds, the small ones first and the large last; a gentle
step recruits perhaps a tenth to a fifth of them.
\textbf{17.} Units are recruited in order of size: at low effort each
new unit adds a small increment, at high effort the remaining units
are large and each adds a big step.
\textbf{18.} The end-plate potential of every fibre is halved; the
fibres whose potential falls below threshold do not fire, so the
\omterm{def:b2:muscle-movement:motorunit}{twitch} and \omterm{def:b2:muscle-movement:motorunit}{tetanus} are smaller, and the jump falls short.
\textbf{19.} $5\times 20 = \qty{0.1}{mol}$ of ATP: six jumps' worth.
\textbf{20.} $20\times 20 = \qty{0.4}{mol}$ of \omterm{prop:b2:muscle-movement:energy}{creatine phosphate}:
about 23 jumps.
\textbf{21.} Ten jumps use $0.17$ mol, not yet exhausting the \omterm{prop:b2:muscle-movement:energy}{creatine
phosphate}; after twenty or so, glycolysis takes over, and lactate
and protons accumulate.
\textbf{22.} \qty{1.8}{kW} of work at \qty{40}{\%} is \qty{4.6}{kW} of
ATP, $0.09$ mol a second, needing $0.015$ mol of oxygen a second ---
\qty{0.34}{L/s}, \qty{21}{L/min}, ten times what the \omterm{def:b2:blood-circulation:blood}{blood} can bring:
the jump is anaerobic by necessity, and the oxygen comes afterward.
\textbf{23.} $0.17$ mol of ATP needs $0.029$ mol of oxygen,
\qty{0.64}{L}; paid over a minute, \qty{640}{mL/min}, two and a half
times the resting consumption.
\textbf{24.} ATP is rephosphorylated from \omterm{prop:b2:muscle-movement:energy}{creatine phosphate} within
milliseconds, so its concentration barely falls; rigor needs ATP
near zero, which a living muscle with any reserve never reaches.
\textbf{25.} 42 strokes available per \omterm{def:b2:cell-differentiation:fibre}{half-sarcomere}, 24 needed;
$F/F_0 = 0.55$ at the jump's speed; $0.017$ mol, \qty{8.7}{g}, of ATP.
\end{solution}
